% NAF 5176, batch 3 — Gallica views 41–58 (the last batch; eighteen views).
% Pass: Opus 5.5 (claude-opus-5-5), 2026-09-22 — first pass, unchecked against the views by a human.
%
% The views. Greyscale microfilm scans, landscape, about 8730 × 6410: each
% view is the volume lying open, a left page and a right page; small leaves
% are mounted on guards and sit anywhere on the opening. View 58 is the back
% board alone (4679 × 6409). In the file, \page{N} opens the left page of
% view N; the right page follows under the same \page, introduced by a \note{}.
%
% What the eighteen views hold. Mathematical working papers in French, on
% games of chance and on combinations, in two hands; no letter, no date, no
% signature, no name of anyone. Nothing on these leaves says who wrote them;
% the catalogue title (Mersenne and Pascal) is collective, and this pass
% attributes nothing.
%
%   v. 41 L    Verso of f. 32 (view 40): the division of the stakes (« partis »)
%              between two players, continued without a heading from the recto —
%              the two-game case, then « En 4 Jeux » (the values 5/32, 5/32, 4/32,
%              2/32 of the successive games), a « Regle gnale » and a remark that
%              equal scores count for nothing. First hand: a fast cursive, black
%              ink, with marginal notes half lost in the gutter shadow.
%   v. 41 R    f. 33r, same hand: « 3 Joueurs En 2 Jeux » (A, B, C: 17/27,
%              5/27, 5/27), « Regle gnale pour tout » (an expectation is the sum
%              of the equally likely outcomes divided by their number), with the
%              proof by an equivalent fair game, and a « Regle » for any number of
%              players.
%   v. 42 L    f. 33v: « Pour les partis des dés » — tables of the 36 throws of
%              two dice and of the ways of making each point with two and with
%              three dice; the odds of throwing a six in 1 to 5 throws of one die
%              (671 against 625 in four throws).
%   v. 42 R    f. 34r: the same continued to 7 throws, with a table; then « A 2
%              dés », the odds of two sixes in 1 to 4 throws, and the remark that
%              in 24 throws there is a little disadvantage and in 25 a little
%              advantage; ends on a question left unanswered.
%   v. 43–46   Blank: 34v (showing the recto through, mirrored), 35r (numbered
%              35; views 43 and 44 are two exposures of the same opening), 35v,
%              36r (a small leaf numbered 36), 36v, 37r (numbered 37), and the
%              guards. Not transcribed.
%   v. 47 L    A blank page; on a strip along its edge, written lengthwise, the
%              dossier title « Geometrie Meslée » (reading of the second word
%              doubtful), not in the hand of the text.
%   v. 47 R    f. 38r, a small half-height leaf in a second hand, round and
%              compact: compound interest on one ducat at 5 per cent for 32 years,
%              by squaring 21/20 five times (the 32nd powers of 21 and 20 written
%              out in full), then 33 and 34 years.
%   v. 48      f. 38v – 39r: « Des Combinaisons Multiples » — six soldiers ranged
%              and given arms and liveries (720, 518400, 373248000; then with
%              repeated arms and liveries, 7776000; then with more kinds than
%              soldiers, 73,156,608,000); the start of a heavily corrected problem
%              first put as bouquets in pots, then as commanders placed in posts.
%   v. 49 L    f. 39v: the commanders continued; then words formed from the eight
%              letters a b c d e i o s with b, c, d never together (36000 « dictions
%              ou anagrammes »), and with two of them never at the start or end
%              (28800). A draft, much struck through.
%   v. 49 R –  f. 40r – 41v: « N.o 2. » / « Hasard » — the yearly draw of five
%   v. 51 L    names at « Gennes » (reading doubtful) and the bets laid on it
%              with bankers: the chances of the better and of the banker when 5,
%              4, 3, 2 or 1 of the five named are drawn out of 100
%              (75,287,520 draws), the banker's odds (about 1 to 3764 3/4 for all
%              five), their sum, and « le fondement de ces raisons », which
%              explains the products and the triangular and tetrahedral numbers
%              used. Stops on view 51 L on a sentence without its end.
%   v. 51 R    f. 42r, a small leaf, same second hand: the number of ancestors in
%              30 generations (2^30 = 1,073,741,824); the ways of moving the
%              eight pawns at chess once each when each may advance one or two
%              squares (10,321,920), or one, two or three (264,539,520), with two
%              tables.
%   v. 52      Left, 42v: blank, the recto showing through mirrored (checked by
%              mirroring the crop). Right, a blank guard. Not transcribed.
%   v. 53–57   Blank guard leaves at the end of the volume. Not transcribed.
%   v. 58      The back board, marbled. Not transcribed.
%
% Across the batch's edges. View 41 L begins in mid-sentence (« ou vn hazard
% egal … »): the partis text begins on view 40, in batch 2. The second hand's
% pieces are self-contained on ff. 38–42, except that « Des Combinaisons
% Multiples » runs from 38v to 39v and « Hasard » from 40r to 41v.
%
% Numbers on the leaves. Two ink series, one number per leaf, at the top right
% of the rectos.
%   (a) A continuous series 32–42, one per leaf, blanks included: 33 (v. 41 R),
%       34 (42 R), 35 (43 R and 44 R), 36 (45 R), 37 (46 R), 38 (47 R),
%       39 (48 R), 40 (49 R), 41 (50 R), 42 (51 R); 32 is on view 40, outside
%       the batch. It is in a thin hand distinct from both texts, stands on blank
%       leaves, runs without a gap, and ends at 42, the catalogue's count of
%       folios, just before the blank guards. It behaves as the library's
%       foliation, and batches 1–2 read the same series (1–32), but it is in
%       ink, not pencil: as for Latin 17859 and Français 9115, no \folio{} is
%       written until a human has looked at the leaves.
%   (b) A smaller, paler, older series beside or below it on the written leaves
%       only, out of order: 24 (f. 33), 25 (f. 34), 27. (f. 38), 30 (f. 39),
%       28 (f. 40), 29. (f. 41), 26 (f. 42). Not a folio; recorded in the notes.
%   Also: « N.o 2. » at the head of f. 40r (a piece number, the writer's or a
%   filer's); the round stamp « BIBLIOTHÈQUE NATIONALE MANUSCRITS R.F. » on
%   ff. 34, 38, 39v, 40, 42. Numbers inside the texts and tables are content.
%
% The hands and the notation. The long s is set as s throughout. Abbreviations
% kept: « gnale » (générale), « comb. », « \&c ». Spellings kept: « escus »,
% « hazard/hazart/hazars », « Commendans », « ancestres », « eschets »,
% « consones ». Fractions are written with a bar and set \frac{}{}; the writers'
% commas between groups of three digits are kept. Tables drawn with pen rules
% are set as array; their layout is described in notes. Old numbers in the
% second hand: the 4 is a closed loop that reads as 9 or 1, and the 5 can look
% like a 3 — every doubtful figure was checked against the arithmetic of the
% page, and where the leaf and the sum disagree the leaf is kept and a note
% says so.
%
% Unsettled. The first line of the Hasard piece (two or three words, perhaps the
% size of the list the five are drawn from); the town (« Gennes »?); several
% corrections of the struck drafts on views 48 R and 49 L, whose final wording
% is not certain; the last line of view 41 L, under a dark stain; the marginal
% notes of view 41 L in the gutter shadow.

\documentclass[11pt,a4paper]{article}
% The preamble shared by every transcription of the mathematicians' manuscripts in Gallica.
%
% It rests on one idea: **uncertainty compounds, it does not resolve.** A
% transcription that smooths over a doubtful word has destroyed the only thing
% it was made for — a reader must be able to tell, at every point, what was on
% the page and what was supplied. The apparatus macros below are therefore the
% critical apparatus, not decoration.
%
% The same commands are recognised by `scripts/render.mjs`, which produces the
% site's reading view, and by `scripts/tei.mjs`. A macro added here must be
% added there too, or rendering fails loudly — which is the wanted behaviour.
%
% Comments here are English, like the rest of the repository. The documents
% this preamble sets are French, and that is deliberate: see the skills.

\ProvidesPackage{germain}[2026/09/15 Transcriptions of mathematicians' manuscripts in Gallica]

\usepackage{amsmath,amssymb,amsthm}
\usepackage{mathrsfs}
% Kept from the parent preamble so the renderer's diagram support stays
% available; these manuscripts draw no commutative diagrams, and nothing here uses it.
\usepackage{tikz-cd}
\usepackage[english,french]{babel}
\usepackage{fontspec}

% --- Greek ------------------------------------------------------------------
% The text face stays fontspec's default, Latin Modern, which has no Greek:
% XeTeX drops every glyph it lacks with a line in the log and nothing on the
% page, and the Greek words the manuscripts quote vanished from the PDFs. So
% the Greek and Greek Extended blocks get a XeTeX character class of their own,
% and entering or leaving a run of it switches the family to CMU Serif — the
% same Computer Modern design, with polytonic Greek — and back. Only the
% family changes: a Greek word in \textit or \textbf keeps its shape and series.
% The transitions to and from every other class carry the tokens that class
% has towards an ordinary letter, so French spacing before « ; » or inside
% « » still holds after a Greek word. No grouping, so no brace or \endgroup
% can come out unbalanced; maths is left alone, since XeTeX inserts nothing
% there. Kept identical in `exercices.sty`.
\makeatletter
\newfontfamily\germainGreekfont{cmun}[
  Extension=.otf, NFSSFamily=germaingreek,
  UprightFont=*rm, ItalicFont=*ti, SlantedFont=*sl,
  BoldFont=*bx, BoldItalicFont=*bi, BoldSlantedFont=*bl]
\newXeTeXintercharclass\germain@greek
\def\germain@greekrange#1#2{%
  \@tempcnta=#1\relax
  \loop
    \XeTeXcharclass\@tempcnta=\germain@greek
    \ifnum\@tempcnta<#2\advance\@tempcnta\@ne\repeat}
\germain@greekrange{"0370}{"03FF}
\germain@greekrange{"1F00}{"1FFF}
\def\germain@greekprev{\rmdefault}
\def\germain@greekin{%
  \edef\germain@greekprev{\f@family}\fontfamily{germaingreek}\selectfont}
\def\germain@greekout{\fontfamily{\germain@greekprev}\selectfont}
\def\germain@greekjoin{%
  \XeTeXinterchartokenstate=\@ne
  \@tempcnta=\z@
  \loop
    \ifnum\@tempcnta=\germain@greek\else
      \XeTeXinterchartoks\germain@greek\@tempcnta=\expandafter{%
        \expandafter\germain@greekout\the\XeTeXinterchartoks\z@\@tempcnta}%
      \XeTeXinterchartoks\@tempcnta\germain@greek=\expandafter{%
        \the\XeTeXinterchartoks\@tempcnta\z@\germain@greekin}%
    \fi
    \ifnum\@tempcnta<4095\advance\@tempcnta\@ne\repeat}
% babel-french sets its own classes, and saves and restores the interchar
% state, each time a language is selected: join again after it does.
\addto\extrasfrench{\germain@greekjoin}
\addto\extrasenglish{\germain@greekjoin}
\AtBeginDocument{\germain@greekjoin}
\makeatother

\usepackage{geometry}
\geometry{a4paper,margin=25mm,marginparwidth=28mm}
\usepackage{marginnote}
\usepackage{xcolor}
\usepackage[normalem]{ulem}
\setlength{\emergencystretch}{3em}
\usepackage{hyperref}
\hypersetup{colorlinks=true,linkcolor=black,urlcolor=black,pdfborder={0 0 0}}

% --- Batch metadata ---------------------------------------------------------
% Taken as they stand by the rendering script, which shows them at the head of
% the reading view. They fix what the file claims to cover. The macro names are
% the parent project's, so that the scripts stay shared; \folder here means the
% volume's slug — `fr-9115` — and \pages counts Gallica views.

\newcommand{\folder}[1]{\gdef\grothFolder{#1}}
\newcommand{\batch}[1]{\gdef\grothBatch{#1}}
\newcommand{\pages}[2]{\gdef\grothFirst{#1}\gdef\grothLast{#2}}
\newcommand{\foldertitle}[1]{\gdef\grothFoldertitle{#1}}

% The holder's own shelfmark, printed as they write it: « Français 9115 ».
\newcommand{\shelfmark}[1]{\gdef\grothShelfmark{#1}}

% The dating, copied verbatim from the holder's catalogue — never inferred
% here. « XIXe siècle » is the BnF's; a pass that dates a leaf from its content
% says so in a \note{}, as its own inference.
\newcommand{\dating}[1]{\gdef\grothDating{#1}}

% The status notice, printed in the title block and repeated as a diagonal
% watermark on every page of the PDF. The manuscripts are in the public
% domain; what this line declares is not a right but a quality — an
% unverified first pass by a machine. The skills write the line; a file
% without it gets no watermark, so leaving it out is visible in the output.
\newcommand{\watermark}[1]{\gdef\grothWatermark{#1}}

\gdef\grothFolder{?}\gdef\grothBatch{}\gdef\grothFirst{?}\gdef\grothLast{?}%
\gdef\grothFoldertitle{}\gdef\grothShelfmark{}\gdef\grothDating{}\gdef\grothWatermark{}

% --- The résumé -------------------------------------------------------------
\newenvironment{resume}
  {\small\setlength{\parskip}{0.45em}}
  {\par\medskip\hrule\medskip}

% --- Keywords ---------------------------------------------------------------
% In English, closing the modernised reading's résumé: the modern vocabulary
% under which to search for what the leaves construct. The manifest extracts
% them to tag the volume — this line is the single source of the tags.
\newcommand{\keywords}[1]{\par\smallskip\noindent
  {\footnotesize\textsc{Keywords} --- \textit{#1}}\par}

% --- The critical apparatus -------------------------------------------------

% The Gallica view number, in the margin. This is the anchor that lets a
% disagreement over a formula be settled by looking at *one* image rather than
% a volume of 750. The site uses it to turn the facsimile as the transcription
% is scrolled. A view is one scanned image: usually one side of a leaf.
\newcommand{\page}[1]{%
  \marginnote{\footnotesize\color{gray!70}\textsf{v.\,#1}}%
  \label{page:#1}\ignorespaces}

% The BnF's pencilled foliation, where the leaf carries it and a pass can read
% it: « 198r ». This is what the literature cites — « ff. 198r–208v » — and
% the one line that turns a citation into a view. Recorded, never inferred:
% a folio a pass cannot read is not written.
\newcommand{\folio}[1]{%
  \marginnote{\footnotesize\color{gray!50}\textsf{f.\,#1}}\ignorespaces}

% The modernised reading's coarser counterpart to \page{}: which views a
% section is drawn from.
\newcommand{\pagerange}[2]{%
  \marginnote{\footnotesize\color{gray!70}\textsf{v.\,#1--#2}}%
  \ignorespaces}

% Illegible: never guessed.
\newcommand{\ill}{\textcolor{red!60!black}{[\,\ldots\,]}}

% A doubtful reading: the word is given, and flagged as uncertain.
\newcommand{\uncertain}[1]{\uline{#1}}

% An editorial addition — a missing letter, an implied word.
\newcommand{\add}[1]{\textcolor{blue!50!black}{[#1]}}

% Struck out by the writer. What was crossed out is part of the document.
\newcommand{\struck}[1]{\sout{#1}}

% The transcriber's note, on the state of the page or the mathematics.
\newcommand{\note}[1]{\par\smallskip{\small\color{gray!60!black}\textit{[#1]}}\par\smallskip}

% A marginal note of hers, distinct from the body of the text.
\newcommand{\marginal}[1]{\par\smallskip\noindent\hspace{1em}%
  {\small\color{gray!60!black}#1}\par\smallskip}

% --- Title ------------------------------------------------------------------
% The title block is French, like the documents themselves.

\AddToHook{shipout/background}{%
  \ifx\grothWatermark\empty\else
    \begin{tikzpicture}[remember picture, overlay]
      \node[rotate=52, scale=3.4, align=center, text=gray!18,
            font=\sffamily\bfseries]
        at (current page.center) {\grothWatermark};
    \end{tikzpicture}%
  \fi}

\AtBeginDocument{%
  \begin{center}
    {\large\bfseries Manuscrits de mathématiciens (Gallica) ---
      \ifx\grothShelfmark\empty\grothFolder\else\grothShelfmark\fi}\\[2pt]
    {\itshape\grothFoldertitle}\\[4pt]
    {\small vues \grothFirst--\grothLast
      \ifx\grothBatch\empty\else\ (lot \grothBatch)\fi}\\[2pt]
    \ifx\grothDating\empty\else
      {\small\color{gray!60!black}Datation du catalogue : \grothDating}\\[2pt]
    \fi
    \ifx\grothWatermark\empty\else
      {\small\bfseries\color{red!45!gray}\grothWatermark}\\[2pt]
    \fi
    {\footnotesize\color{gray!60!black}Transcription établie d'après les images IIIF de Gallica.
     Source gallica.bnf.fr / Bibliothèque nationale de France.\\
     Les lectures incertaines sont soulignées, les illisibles notées [\,\ldots\,],
     les ajouts entre crochets ; \textsf{v.} numérote les vues de Gallica,
     \textsf{f.} la foliotation au crayon de la BnF.}
  \end{center}
  \vspace{1em}}

\endinput


\folder{naf-5176}
\shelfmark{NAF 5176}
\batch{3}
\pages{41}{58}
\dating{1601-1700}
\watermark{Lecture automatique\\non vérifiée}
\foldertitle{Papiers du Père MERSENNE et de Blaise PASCAL.}

\begin{document}

\page{41}
\note{Verso du feuillet 32 (vue 40) : le texte continue sans titre la page commencée au recto. Écriture cursive rapide, à l'encre noire ; la marge de gauche est prise dans l'ombre de la reliure, qui cache le début des lignes des notes marginales.}

ou vn hazard egal de gagner ou de perdre le jeu suivant. S'il le gagne Il gagne tout — S'il le perd Il n'aura plus que deux jeux a vn. donc B en aura vn, Et A en a 2 — donc A \struck{a} \uncertain{aura} les $\frac{3}{4}$ dans le tout qui sont $\frac{12}{16}$ et B $\frac{4}{16}$ — donc Il faut dire qu'A ayant 2 jeux a \struck{\ill{}} rien avait egallement hazard pour avoir 16 escus ou pour en avoir 12 — donc les 12 luy sont seurs, Et il faut partager egallement le surplus, qui est 4. dont la moitié est 2. qui joints a 12 font 14. — donc il avait 14 escus et B seulement 2. — ou autrement Il n'y a qu'a partager les esperances de 16, et 12. dont la moitié est 14. ce qui est \uncertain{la} mesme chose que j'ay dit
\note{« aura » est écrit au-dessus de la ligne, sur un « a » biffé. Devant « rien », un chiffre biffé et noyé d'encre ; « rien » est écrit au-dessus. Les derniers mots, « ce qui est la mesme chose que j'ay dit », sont serrés sous la ligne, au début de la suivante.}

Pareillement suivant \uncertain{comme s'ensuit} — quand A a vn jeu a rien il a vn egal hazard d'avoir le jeu suivant ou de le perdre. S'il le gagne il luy reuient 14 escus, S'il le perd il en aura encor \struck{12} 8 — donc \struck{\ill{}} Il est en hazard egal d'en avoir \struck{\ill{}} 8 ou 14 — donc partageant ces deux esperances de \uncertain{semblable} la moitié est 11 \struck{\ill{}} — donc ayant vn jeu a rien il luy reuient 11 escus et a l'autre seulement 5.
\note{le 8 est écrit au-dessus d'un nombre biffé, qui semble être 12.}

\subsection*{En 4 Jeux}

A en a 2 a rien. combien luy reuient il?

(\uncertain{Posons} le cas qu'il ne luy en manquast qu'vn Et a B, 2 — A aurait $\frac{12}{16}$ du jeu. Car c'est toujours comme sy on jouait \uncertain{vn} \ill{} coup Et qu'on en fust vn a rien) donc A aurait 3 jeux Et B. 2 — donc Il faut dire, quand A a 3 jeux a 1 il a aussy hazard pour avoir $\frac{12}{16}$ (sy B fait le suivant) que pour avoir $\frac{16}{16}$ sy A le gagne — donc il luy reuient — $\frac{14}{16}$. Pareillement quand A a 3 jeux a rien il a aussy hazard pour $\frac{14}{16}$ (sy B fait le suivant) que pour $\frac{16}{16}$ (sy A le gagne) — donc A ayant 3 jeux a rien a $\frac{15}{16}$. — ainsy sy j'en ay 2 a rien il faut dire — J'ay hazard egal pour avoir 3 a rien ou pour avoir 2 a I. Sy J'ay 3 a rien J'auray $\frac{15}{16}$. Sy J'ay 2 a I, Il luy faudra encor 3 jeux Et a moy 2. S'il luy faut 3 jeux et a moy 2, c'est comme sy \uncertain{j'en} avais vn a rien de 3. \struck{donc} auquel cas j'aurais $\frac{11}{16}$. donc ayant 2 a 0, J'ay hazard egal pour $\frac{15}{16}$ ou pour $\frac{11}{16}$. donc J'ay la moitié des 2 hazards qui sont $\frac{13}{16}$.
\note{après « vn a rien », au-dessus de la parenthèse fermante, un petit « 16 » isolé. Dans la première parenthèse, « sy B fait le suivant » est écrit au-dessus d'un mot biffé.}
\marginal{\ill{} Autrement — c'est qu'on a \ill{} egal pour avoir $\frac{16}{16}$ ou \ill{}, c'est avoir \ill{}}
\note{note de la marge de gauche, en face de « quand A a 3 jeux a 1 » ; le début des lignes est pris dans l'ombre de la reliure, et les deux dernières fractions sont noyées d'encre : on y distingue seulement le dénominateur 16.}
\marginal{car c'est comme sy on \ill{} en 3 de quoy l'on \uncertain{suppose} \ill{} qui est clair, parce \uncertain{que} d'vn jeu nous serons vn \ill{} a partager le tout pour vn 3.}
\note{note de la marge de gauche, en face de « jeux Et a moy 2 … auquel cas j'aurais » ; même ombre de la reliure au début des lignes.}

plus court. Sy J'ay 2 a I. J'ay egal hazard pour $\frac{14}{16}$ ou pour $\frac{8}{16}$ — donc J'ay $\frac{11}{16}$.

donc sy J'ay 2 a 0 J'ay hazard egal pour $\frac{11}{16}$, ou pour \struck{4} $\frac{15}{16}$ — donc J'ay $\frac{13}{16}$.

donc sy J'ay 1 a 0 J'ay hazard egal pour 2 a 0, ou pour 1 a vn — c'est a dire pour $\frac{13}{16}$ ou pour $\frac{8}{16}$

donc J'ay $\frac{21}{32}$.

Tellement que Jouant en 4 Jeux le premier vaut $\frac{5}{32}$ qui est $\frac{21}{32}$ au total

le 2 vaut encor — — $\frac{5}{32}$ qui est $\frac{26}{32}$ au total

le 3 vaut — — — $\frac{4}{32}$ qui est $\frac{30}{32}$ au total

le 4 vaut — — — $\frac{2}{32}$ qui est en tout $\frac{32}{32}$

Toute cela supposé qu'vn seul les gagne tous 4 de suitte.
\note{le tableau est en quatre lignes, les « vaut » reliés aux fractions par des pointillés, les sommes tirées à droite et séparées par des traits. Dans la première ligne le numérateur de « $\frac{21}{32}$ au total » est en partie noyé d'encre.}
\marginal{N\textsuperscript{a} A vaut mieux avoir 2 jeux a 0 de \uncertain{4} qu'vn a 0 de deux. \ill{} de $\frac{1}{16}$ davantage}
\note{note de la marge de gauche, séparée du texte par un trait et une accolade, en face des lignes « le 4 vaut » et « Toute cela supposé ». Les débuts de ligne sont dans l'ombre de la reliure.}

\subsection*{Regle gnale}

Pour \uncertain{venir} a la connoissance de la valeur du premier jeu Il faut commancer par le plus pres de la fin c'est a dire par le cas auquel Il faut vn jeu a l'vn, Et 2 a l'autre Et ce cas la vaut toujours $\frac{3}{4}$ de tout l'argent a celuy a qui il ne faut qu'vn jeu.
\note{« de tout l'argent » est écrit au-dessus de la ligne, après la fraction.}

\subsection*{\uncertain{Autre}}
\note{titre centré sous la règle ; le mot, d'une grande initiale et de quatre lettres, a la forme du « Autre » biffé plus bas.}

La valeur des jeux change a mesure que chacun en fait alternativement, car des qu'on \uncertain{est} egaux \struck{qu'ils sont egaux} du nombre de jeux, toute la somme se joue dans les jeux qui restent. par exemple quand on joue en 4 jeux, le 1\textsuperscript{er}. vaut naturellement $\frac{5}{32}$. mais des qu'on est vn a vn, cela se compte pour rien Et c'est comme qui jouait tout en trois. auquel cas le premier d'apres vaut $\frac{3}{16}$. que sy on vient encor a estre egaux — le tout se joue en 2 auquel cas le premier vaudra $\frac{3}{4}$.
\note{« qu'on est egaux » est écrit au-dessus de la ligne, sur « qu'ils sont egaux » biffé d'un trait ; le mot du milieu de l'ajout paraît lui-même repassé.}

\struck{Autre} plus claire

Des qu'on est egaux en jeux, ces jeux la sont comptés pour rien Et la partie se joue au nombre des jeux suivans.
\note{« Autre » est biffé d'une croix dont un bras descend en diagonale à travers les deux lignes suivantes ; on ne peut dire si ce trait biffe aussi le paragraphe, qui est donné ici sans biffure.}

\ill{} vn jeu \ill{} et B le perd, \ill{} A le \ill{}
\note{dernière ligne de la page, dans la bande sombre d'une tache qui couvre le pied du feuillet ; seuls quelques mots se distinguent, même en poussant le contraste.}

\note{Page de droite de la vue : recto du feuillet suivant, même main, écrite plus posément. En haut à droite, deux nombres à l'encre : « 24 », petit et fin, et au-dessous « 33 », plus grand et plus appuyé ; le second est celui de la série continue (un nombre par feuillet, blancs compris) qui court de 32 à 42 sur ce lot et que les trois lots lisent comme la foliation probable, non donnée comme folio (voir l'en-tête) ; le premier appartient à une numérotation plus ancienne, qui n'est pas dans l'ordre des feuillets.}

\subsection*{3 Joueurs En 2 Jeux}

A.B.C.

A a 1 Jeu, B, 1, Et C 0. combien revient il a chacun?

Il est egal a \struck{A d'avoir} tous les 3 d'avoir le jeu suivant (car chacun a 2 hazards contre soy) Et par consequent en quittant la partie chacun a $\frac{1}{3}$ de la valeur de ce jeu.

donc Il est egal a A d'avoir ou $\frac{27}{27}$ (s'il gagne ce jeu) ou $\frac{9}{27}$ (sy C gaigne) ou 0 sy B le gagne. donc partageant ces trois esperances par tiers parce qu'ils sont trois joueurs, Il luy revient $\frac{12}{27}$, Et a B $\frac{12}{27}$ Et a C $\frac{3}{27}$

A a vn Jeu Et B 0, et C 0. combien chacun doit il avoir en quittant?

Il est egalement hazardeux pour A d'avoir $\frac{27}{27}$ (s'il gaigne le jeu suivant) ou $\frac{12}{27}$ (sy B le gaigne) ou $\frac{12}{27}$ sy C le gagne — partageant ces trois hazards par 3, Il revient $\frac{17}{27}$ au \uncertain{quotient} sur $\frac{51}{27}$ — donc A aura en ce cas $\frac{17}{27}$, B $\frac{5}{27}$, et C $\frac{5}{27}$.

autrement on peut dire, A au pis aller aura toujours $\frac{12}{27}$ soit que B, ou C gaignent ce jeu. donc il a ces $\frac{12}{27}$ asseurés, Et comme il peut gagner ce jeu aussy bien qu'eux Il y a son tiers comme eux au Restant. or le Restant sont $\frac{15}{27}$ — donc Il luy en revient $\frac{5}{27}$ qui joins a ses 12, font $\frac{17}{27}$.

\subsection*{Regle gnale pour tout}

Sy J'ay vne egale attente Et facilité a obtenir A ou B ou C. mon esperance \struck{attente} vaut le tiers de \struck{A+} A+B+C. — Sy vne egale facilité a obtenir A, B, C, ou D, mon esperance vaut le quart de A+B+C+D. Et ainsy a l'infiny. c'est a dire qu'on a toujours le quotient produit de l'addition des sommes qu'on espere a jeu egal. \struck{c'est} divisées par le nombre des attentes
\note{« esperance » est écrit au-dessus de « attente » biffé ; « attentes ou des cas » est écrit au-dessus de « des sommes ». Le mot « produit » est repassé à l'encre plus grasse.}

Ainsy sy J'ay \struck{\ill{}} 3 cas ou 3 coups pour avoir 12, Et \uncertain{3} coups pour avoir 8 cela me vaut autant que 10. Et cela revient a la premiere Regle. car ayant 10 Il est aisé de voir que je puis revenir a la mesme esperance Jouant a jeu egal. car Jouant contre cinq amis qui mettront chacun 10, \struck{a condition} en convenant avec 2 des cinq que celuy qui gaignera me donnera 12, Et avec 3 autres que celuy qui gaignera me donnera 8, ou que sy je gaigne je donneray aux 2 premiers a chacun 12, Et a chacun des 3 autres 8. — Il est constant que J'ay 3 attentes pour 12, Et 3 pour 8. car sy je gaigne tout J'ay 60, dont donnant \struck{a} 2 fois 12 \struck{aux a} aux vns, Et 3 fois 8 aux autres, Il me reste 12, ou sy l'vn des 2 premiers gagnent il me revient encor 12, sy vn des 3 autres Il me revient 8.
\note{le chiffre devant « coups pour avoir 8 » est surchargé : le trait visible se lit 4, écrit sur ou sous un autre chiffre ; la suite (« 3 attentes pour 12, Et 3 pour 8 ») donne 3.}

Pour rapporter donc l'exemple des 3 joueurs cy dessus a cette Regle, dans le cas que A a 1 jeu de 2. B 0, et C 0. Il faut dire A a vne egale esperance pour 27, pour 12, Et pour 12, c'est a dire vne pour 27, Et 2 pour 12. or ces 3 esperances assemblées font 51, qui divisez par 3 donnent 17 — donc Il a 17. La preuve est que sy A Joue 17 contre B 17, Et C 17 a condition que sy A gagne Il donnera a B 27, Et a C. 12 que sy B gagne il donnera a A 27, Et que sy C gagne Il donnera a A 12. le jeu sera egal et Juste, Et cependant A revient a ses trois esperances. car s'il gagne Il donnera 27 a B, Et 12 a C, Et Il luy restera encor 12. Sy B gagne A en \uncertain{receura} 27, Et sy C gagne A en \uncertain{receura} 12. de sorte qu'il a toujours comme auparavant vne esperance pour 27, Et 2 pour 12.

\subsection*{Regle}

De quelque nombre de Joueurs qu'on soit dont les vns ayent plus de jeux que les autres on doit voir quelle part revient a \struck{chacun} l'vn d'eux supposant qu'il gagne le jeu suivant, ou supposant que ce soient les vns ou les autres qui \uncertain{restent}. Et ayant trouvé tout ce qui luy revient en chaque cas, Il faut diviser \struck{\ill{}} le nombre de toutes ses portions assemblées par le nombre des joueurs, le quotient donne ce qui luy appartient en cet estat. — Les exemples precedens sont calculés la dessus. Et pour trouver le calcul il faut toujours commancer par le cas le plus simple. de la on vient aux autres.
\note{« l'vn d'eux » est écrit au-dessus de « chacun » biffé.}

\page{42}
\note{Page de gauche de la vue, sans numéro ; même main que les pages précédentes, plus posée. Le coin supérieur gauche du feuillet est coupé.}

\subsection*{Pour les partis des dés}

Vn dé estant cubique Il a 6 faces Et ainsy il peut amener 6 coups differens

2 dés amenent 6 fois 6 qui sont 36 comme on peut voir par cette petite table ou on voit que la combinaison En marquant la difference des coups du Dé A et du dé B

A peut faire — pendant que B fait
\[
\begin{array}{c|cccccc}
1 & 1 & 2 & 3 & 4 & 5 & 6 \\
2 & 2 & 3 & 4 & 5 & 6 & 1 \\
3 & 3 & 4 & 5 & 6 & 1 & 2 \\
4 & 4 & 5 & 6 & 1 & 2 & 3 \\
5 & 5 & 6 & 1 & 2 & 3 & 4 \\
6 & 6 & 1 & 2 & 3 & 4 & 5
\end{array}
\]
\note{table à droite du texte, en cases tracées à la plume ; « A peut faire » est écrit au-dessus de la colonne des chiffres de gauche, en partie enfermé dans un trait.}

3 dés font 6 fois 36 coups qui sont 216

4 dés font 6 fois 216. coups qui sont 1296 Et ainsy du reste dont Il n'est pas besoin de faire les tables.
\note{« fois » est écrit au-dessus de la ligne, entre « 6 » et « 216 ».}

2 dés \uncertain{ne peuvent} faire 2 poincts que par vn coup — qui est 1 et 1

Ils font 3 par 2 coups scavoir A 1, Et B 2, ou A 2, Et B 1.

mais cela se connaist aisement pour peu qu'on y pense. En voicy la table

a 2 dés — points — le font en
\[
\begin{array}{ccl}
2 \text{ et } 12 & \cdots & 1 \\
3 \text{ et } 11 & \cdots & 2 \\
4 \text{ et } 10 & \cdots & 3 \\
5 \text{ et } 9 & \cdots & 4 \\
6 \text{ et } 8 & \cdots & 5 \\
7 & \cdots & 6 \\
\hline
 & & 21 \\
 & \text{Et} & 15 \\
\hline
 & & 36
\end{array}
\]
a 3 dés
\[
\begin{array}{ccl}
3 \text{ et } 18 & \cdots & 1 \\
4 \text{ et } 17 & \cdots & 3 \\
5 \text{ et } 16 & \cdots & 6 \\
6 \text{ et } 15 & \cdots & 10 \\
7 \text{ et } 14 & \cdots & 15 \\
8 \text{ et } 13 & \cdots & 21 \\
9 \text{ et } 12 & \cdots & 25 \\
10 \text{ et } 11 & \cdots & 27 \\
\hline
 & & 108 \\
 & \text{et} & 108 \\
\hline
 & & 216
\end{array}
\]
\note{les deux tables sont côte à côte sur la page, celle des 3 dés à gauche ; elles sont données ici l'une après l'autre. Dans la table des 2 dés, une accolade réunit les nombres de façons, avec « façons » à droite, et en marge de droite « a 2 dés 36 coups ». Dans celle des 3 dés, les mots « points se peuvent faire en » sont écrits en travers des pointillés, « façons. » à droite du 10 ; en marge de gauche « a 3 dés 216 coups », et sous le total 216, le mot « partis ».}

on propose de faire vn 6 avec vn dé; En combien de coups le peut on prendre?

Sy on le prend en vn coup. il faut seulement jouer 1 contre 5, Et sy on devait quitter sans jouer. celuy qui me propose le party aurait $\frac{5}{6}$ Et moy $\frac{1}{6}$. car suivant les Regles cy devant J'ay 5 esperances pour 0 Et seulement 1 pour 6. — donc partageant ces esperances par 6 Il me revient $\frac{1}{6}$. — Et qu'il les faille partager par 6 Il est evident encor que c'est comme sy nous estions 6 joueurs \uncertain{contre 5 desquels j'aurois mis} \struck{qui \ill{} chacun} d'amener 6 du premier coup. or qu'ils soient 5 gageant chacun 1, ou qu'il n'y ait qu'vn gageant 5 c'est la mesme chose.
\note{« contre 5 desquels j'aurois mis » est écrit au-dessus de la ligne, sur trois ou quatre mots biffés d'un trait.}

Sy Je le prends en 2. J'ay desja pour mon 1\textsuperscript{er}. coup $\frac{1}{6}$, Et pour le second $\frac{1}{6}$ de ce qui reste ce qui est comme $\frac{11}{36}$. car posons que nous jouions sur 36 escus. le 1\textsuperscript{er}. coup m'en donne 6, Et a l'autre 30. le 2\textsuperscript{e} coup me vaut pareillement $\frac{1}{6}$ de 30 qui est 5, Et joignant 6, a 5 J'ay $\frac{11}{36}$, Et il en reste a l'autre $\frac{25}{36}$. de sorte qu'a prendre en 2 coups Je puis jouer 11 contre 25
\note{sous 11 et sous 25, le trait de fraction et le dénominateur 36 sont biffés.}

Au 3\textsuperscript{e}. coup J'ay encor la 6\textsuperscript{e} partie de $\frac{25}{36}$ qui est $\frac{25}{216}$ qui joint a $\frac{11}{36}$ font $\frac{91}{216}$. donc a le prendre en 3 Je puis jouer 91, contre \struck{216} 125. Car Il reste a l'autre $\frac{125}{216}$

Au 4\textsuperscript{e}. coup. J'ay encor la 6\textsuperscript{e} partie de $\frac{125}{216}$ qui est $\frac{125}{1296}$ qui joint a $\frac{91}{216}$ que J'avais font en tout $\frac{671}{1296}$ donc Il reste a l'autre $\frac{625}{1296}$ — donc Je puis jouer 671 contre 625. qui est plus qu'vn contre vn* donc a prendre a amener vn 6 en 4 coups avec vn dé Il y a vn peu d'avantage.
\note{« a prendre » est écrit en lettres plus grandes que le reste de la ligne. L'astérisque après « contre vn » est de l'auteur ; aucun renvoi correspondant n'a été vu sur la page.}

Au 5\textsuperscript{e} coup J'ay encor la 6\textsuperscript{e} partie de $\frac{625}{1296}$ qui est $\frac{625}{7776}$ Et joint a $\frac{671}{1296}$ que J'avais me fait en tout $\frac{4651}{7776}$ — Et reste a l'autre seulement $\frac{3125}{7776}$ — donc Je puis jouer en prenant en 5 coups 4651 contre 3125 — qui est vn peu moins que 3 contre 2.

\note{Page de droite de la vue : le calcul continue, même main. En haut, au milieu de la marge supérieure, « 34 » à l'encre, de la série continue (voir l'en-tête) ; plus bas à droite, sous la deuxième ligne, « 25 », plus petit et plus fin, de l'ancienne numérotation. Le cachet rond « BIBLIOTHÈQUE NATIONALE MANUSCRITS R.F. » est apposé à droite, vers le bas de la page.}

Au 6\textsuperscript{e}. coup J'ay encor la 6\textsuperscript{e}. partie de $\frac{3125}{7776}$ qui est $\frac{3125}{46656}$ Et joints a $\frac{4651}{7776}$ que J'avais me font en tout $\frac{31031}{46656}$ Et reste a l'autre $\frac{15625}{46656}$ — donc Je puis jouer en 6 coups 31031 contre 15625, qui est vn peu moins que 2 contre 1.

Au 7\textsuperscript{e} coup. J'ay encor la 6\textsuperscript{e} partie de $\frac{15625}{46656}$ qui est $\frac{15625}{279936}$ Et joint a $\frac{31031}{46656}$ que J'avais me font en tout $\frac{201811}{279936}$. Et reste a l'autre $\frac{78125}{279936}$ — donc Je puis jouer (a prendre En 7 coups) 201811 contre 78125. qui est vn peu plus de 5 contre 2.

On peut ainsy calculer tous les autres de tant de coups qu'on voudra. Et on voit par la qu'on ne peut Jamais trouver vn nombre de coups dans ce party la pour jouer jeu egal.

Voicy la table des 7 \struck{\ill{}} coups cy dessus.
\note{la légende est encadrée d'un trait ; un ou deux mots biffés après « 7 ». À gauche de la table, sous la légende : « qui prend a amener 6 en ».}

\[
\begin{array}{rrcrl}
1 & 1 & \text{contre} & 5 & \\
2 & 11 & & 25 & \text{vn peu moins que} \\
3 & 91 & & 125 & \text{vn peu moins que 3 contre 4} \\
4 & 671 & & 625 & \text{vn peu plus que 1 contre 1} \\
5 & 4651 & & 3125 & \text{vn peu moins que 3 contre 2} \\
6 & 31031 & & 15625 & \text{vn peu moins que 2 contre 1} \\
7 & 201811 & & 78125 & \text{vn peu plus que 5 contre 2}
\end{array}
\]
\note{la table relie chaque nombre de coups aux deux mises par des traits. Entre les rangs 4 et 7, en travers des traits de gauche, sur trois lignes : « Coups avec / vn Dé \uncertain{pour} / jouer ». Le mot « Contre » est écrit une seule fois, en grand, au rang 4 ; au-dessus, dans la même colonne, une pile de chiffres surchargés et biffés. Au bout du rang 1, un mot noyé d'encre, puis « 1 contre \ill{} » ; au bout du rang 2, après « vn peu moins que », les mots sont effacés d'une tache d'encre.}

\subsection*{A 2 dés}

J'entreprens de faire 2 six. En combien de coups le dois je faire

Au 1\textsuperscript{er}. coup J'ay $\frac{1}{36}$ Et mon joueur $\frac{35}{36}$ — donc En vn coup Je dois seulement jouer 1 contre 35.

Au 2\textsuperscript{e} coup J'ay encor $\frac{1}{36}$ des 35 restantes qui vaut $\frac{35}{1296}$ Et joint a $\frac{1}{36}$ que J'avais me donnent en tout $\frac{71}{1296}$, Et restent a \struck{l'autre} joueur $\frac{1225}{1296}$ — donc En 2 coups Je dois seulement jouer 71, contre 1225.

Au 3\textsuperscript{e}. coup Il me revient encor $\frac{1}{36}$ des $\frac{1225}{1296}$ restantes. qui vaut \struck{\ill{}} $\frac{1225}{46656}$ Et \struck{reste} \struck{a l'autre} joint aux $\frac{71}{1296}$ que J'avais J'auray en tout $\frac{3781}{46656}$ Et reste a l'autre $\frac{42875}{46656}$

donc En 3 coups Je ne dois jouer que 3781 contre 42875.
\note{devant $\frac{1225}{46656}$, une première fraction est biffée : le numérateur 1225 se lit, le dénominateur est surchargé. La ligne « donc En 3 coups » est soulignée. Dans « $\frac{3781}{46656}$ » de la ligne précédente, le dernier chiffre du numérateur est tracé à part, un peu plus haut.}

Au 4 Il me revient encor $\frac{1}{36}$ des $\frac{42875}{46656}$ restantes, qui fait $\frac{42875}{1679616}$ Et joint aux $\frac{3781}{46656}$ que J'avais me revient en tout a ce quatriesme coup $\frac{178991}{1679616}$ Et reste a l'autre joueur $\frac{1500625}{1679616}$

donc En 4 Coups Je ne dois jouer que 178991 contre 1500625.
\note{la ligne est soulignée.}

En poursuivant les nombres deviennent trop grands pour en faire icy le calcul. c'est assez de se souvenir qu'il n'y a point de certain nombre de coups auquel on puisse jouer jeu egal entreprenant a faire deux six \uncertain{a} deux dés. mais seulement qu'a le prendre En 24 coups on a vn peu de desavantage, Et qu'en 25 coups on a vn peu \struck{plus} d'avantage

Avec combien de dés entreprendre de faire du premier coup deux six
\note{la question reste sans réponse ; le bas de la page est blanc.}

\page{47}
\note{La page de gauche de la vue est blanche. Sur la bande de papier qui longe son bord droit, écrit à l'encre dans le sens de la hauteur, un titre de dossier : « Geometrie \uncertain{Meslée} ». Il n'est pas de la main du texte qui suit.}

\note{Page de droite : un petit feuillet, de la moitié de la hauteur des précédents, d'une autre main, ronde et serrée, à l'encre noire. En haut à droite, « 38 » de la série continue (voir l'en-tête), et à sa droite « 27. », plus fin, de l'ancienne numérotation. Le titre de la page suivante, « Des Combinaisons Multiples », transparaît à l'envers en haut du feuillet. Cachet rond « BIBLIOTHÈQUE NATIONALE MANUSCRITS R.F. » à droite, au milieu de la page.}

Vn homme met vn ducat a la banque a Multiplico pour 32. ans a la charge d'en avoir les interests, \& interest d'interest a raison de 5. pour 100. On demande a combien se montera le principal, \& les interests au bout de ce \uncertain{temps}.

L'interest est $\frac{1}{20}$ par an. Donc au bout de l'an le principal avec l'interest se montera a $\frac{21}{20}$ de ducat.

Pour les années suivantes il faut prendre les puissances du Numerateur \& du Denominateur de cette fraction \& en faire vne fraction, \& le nombre des années sera l'exposant de ces puissances. Il faudra donc prendre la 32\textsuperscript{e}. puissance de 21. \& de 20. pour le principal, \& les interests de 32. ans. J'ay pris 32. pour la commodité de la multiplication, a cause qu'il n'y aura qu'a prendre le quarré des nombres precedens. Ainsy
\note{« ces » est écrit au-dessus de la ligne, entre « de » et « puissances ».}

Pour 2. ans on aura $\frac{441}{400}$ de ducat.

Pour 4. ans $\frac{194481}{160000}$ qui se trouve prenant le quarré de $\frac{441}{400}$

Pour 8. ans $\frac{37822859361}{25600000000}$

Pour 16. ans $\frac{1430568690241985328321}{655360000000000000000}$

Enfin pour 32. ans \struck{\ill{}} on aura $\frac{2046526777500669368329342638102622164679041}{429496729600000000000000000000000000000000}$
\note{le 4 de cette main, fermé en boucle, se confond avec un 1 ou un 9 : les dénominateurs des deux premières lignes se lisent d'abord 100, et le même tracé fait les 4 de 441. Les rangées de zéros des dénominateurs de 16 et de 32 ans n'ont pu être comptées une à une avec certitude ; on les donne au nombre qu'appelle le numérateur, qui se lit chiffre à chiffre. Les fractions de 8 et de 16 ans sont sur la même ligne, séparées par un trait vertical.}

qui sont 4. ducats \struck{\ill{}} $\frac{3}{4}$ \uncertain{peu plus}.

Le profit sera donc 3. ducats $\frac{3}{4}$ \& vn peu plus. \& ainsy il sera presque quadruple du principal.
\note{la fin de la ligne précédente est biffée d'un trait épais sur toute sa longueur, avec une fraction de dénominateur 4 ; « $\frac{3}{4}$ peu plus » est écrit après la biffure.}

Si on voulait continuer a calculer l'interest pour quelques autres années, il faudroit multiplier le numerateur par 21. \& le denominateur par 20. mais pour eviter la difficulté de ces grands nombres, il faut retrancher vne partie de la fraction qui n'apporte point de profit considerable, \& prendre seulement $\frac{20465267775}{4294967296}$ Donc pour 33. ans on auroit $\frac{429770623275}{85899345920}$ qui est vn peu plus de 5. ducats, qui \uncertain{feroient} 4. ducats de profit. Et si on les laissoit encor vne autre année, on auroit $\frac{9025183088775}{1717986918400}$ qui sont 5. ducats $\frac{1}{4}$ peu plus; c'est donc 4. ducats $\frac{1}{4}$ de \struck{pro} profit.
\note{« Et si on les laissoit » est d'une encre plus pâle que le reste de la ligne.}

\page{48}
\note{Page de gauche : verso du petit feuillet de la vue 47, même main ronde et serrée. L'encre du recto transparaît fortement, à l'envers, entre les lignes (les grands nombres des puissances de 21 et de 20 s'y voient en miroir) ; elle n'est pas transcrite ici.}

\subsection*{Des Combinaisons Multiples.}

On a consideré 2. sortes de combinaison, l'vne d'ordre, \& l'autre de varieté ou \uncertain{changement}, \struck{\ill{}} multiplier. \uncertain{En voicy des exemples}.
\note{la deuxième ligne est barrée en grande partie d'un trait qui court sous les mots ; seuls « varieté ou changement » au début et « multiplier » à la fin se lisent sûrement, et l'on ne peut dire exactement où commence la biffure.}

Six soldats se peuvent ranger en 720. façons, \uncertain{que je suppose} par exemple 6. soldats, a qui je donneray 6. sortes d'armes, par ce que les armes se peuvent aussi combiner en 720. façons, il faudra prendre le quarré de 720. \& on aura 518400. façons de les ranger \& de les armer; Mais si outre cela on leur donnoit 6. sortes de livrées il faudra prendre le cube de 720. pour la varieté. Donc on les pourra ranger \struck{\&} \struck{\ill{}} leur donner divers armes \& livrées, \struck{\ill{}} \ill{} qui se fera en 373248000.
\note{« Donc » en tête de la dernière ligne est lui-même biffé d'un trait.}

Mais si on n'avoit pas autant de sortes d'armes \add{\& de livrées} que de soldats par exemple si on n'avoit que \struck{quatre} 4. sortes d'armes \& que de l'vne des sortes on \uncertain{eust} 3. \uncertain{comme} si on avoit 3. \uncertain{espées} vn mousquet vne pique \& vne halebarde, \& qu'on \uncertain{n'eust} que 3. sortes de livrées \uncertain{sçavoir} 2. de chasque sorte. Il faudra prendre pour les armes la combinaison de 6. choses desquelles il y en auroit 3. semblables; or on a fait voir que pour avoir cette combinaison il faut diviser la combinaison de 6. sçavoir 720. par 6. qui est la combinaison des 3. choses semblables \struck{\ill{}} \& on aura 120. pour les diverses façons \struck{donc} on peut armer les soldats. On multipliera donc 720. par 120. \& on aura 86400. manieres de ranger \& d'armer les soldats. Pour celles livrées par ce qu'il y en a 2. de chasque sorte \& de 3. sortes en tout, il faudra diviser 720. par la combinaison de 2. choses repetée 3. fois qui est 8 par ce que 2. multiplié par 2. donne 4. qui estant encor multiplié par 2. donne 8. divisant donc 720. par 8. on aura 90. par lequel il faudra multiplier le produit qu'on vient d'avoir qui est 86400. \& on aura en tout 7776000 manieres d'arranger ces soldats avec ces armes \& livrées differentes.
\note{« \& de livrées » est écrit au-dessus de la ligne. Un trait vertical sépare « d'armer les soldats. » de « Pour celles livrées ». Le 4 de cette main, fermé en boucle, se lit d'abord 9.}

Que si au contraire on avoit de plus de sortes d'armes \& de livrées qu'il n'y a de soldats; par exemple si on avoit 7. sortes

\note{Page de droite : recto du feuillet suivant, même main ; le texte continue sans rupture. En haut à droite, « 39 » de la série continue, et à sa droite « 30 », plus fin et plus pâle, de l'ancienne numérotation.}

d'armes, \& 8. sortes de livrées \struck{dont} qu'on voulust donner aux soldats en toutes les manieres possibles; on se serviroit de la combinaison de varieté, \struck{par ce qu'il y a plus d'armes \& de livrées que de soldats}, voicy comme se fait cette combinaison.

Pour les armes par ce qu'il y en a de 7. sortes, mais qu'on n'en prend que 6. a \struck{chaque} fois, on multipliera 7. par les nombres moindres, jusques a ce qu'on ayt 6. nombres sçavoir jusques a 1. on multipliera donc 7. par 6. le produit 42. par 5. puis ce produit par 4. 3. \& 2. on aura 5040. qui est le mesme nombre que celuy \struck{\ill{}} de la combinaison d'ordre de 7. choses, par ce que 1. ne multiplie point.
\note{dans « le produit 42. », le 4 a la forme d'un 1 ; le calcul donne 42.}

Pour les livrées, par ce qu'il y a 8. choses \& qu'on ne se sert que de 6. il faudra prendre le \uncertain{changement} de 6. choses \uncertain{prises} dans 8. qui se trouve multipliant l'vn par l'autre 6. nombres commençans par 8. en diminuant, sçavoir 3. 4. 5. 6. 7. 8. dont le produit est 20160. y compris l'ordre.

Il faudra donc multiplier 720. qui monstre la quantité des façons dont on peut ranger 6. soldats, par 5040. qui est la varieté des diverses armes, \& le produit 3628800 (qui contient l'ordre des soldats \& des diverses manieres dont on les peut armer) par 20160. qui est la varieté des livrées, \& on aura 73,156,608,000. qui monstre toutes les façons d'arranger les soldats \& de leur donner diverses armes \& livrées.
\note{le dernier produit est écrit avec des virgules entre les tranches de trois chiffres, comme ici.}

\ill{} \struck{\ill{}} combinaison de varieté se peut aussy multiplier par vne autre combinaison de varieté.

On a 6. \struck{pots \& 6. bouquets differens} \add{places a pourvoir de Commendans, mais} \ill{} on \uncertain{n'en veut mettre} 3. de \struck{ces} bouquets \add{l'vn apres l'autre dans 3. de six} \struck{dans 3. de ces} pots en toutes les façons possibles. On n'a \struck{\ill{}} point \uncertain{égard} a l'ordre \struck{\ill{}} pour mettre \struck{dans} les bouquets, car il n'importe pas lequel on mette le premier dans les pots, mais cet ordre est dans \add{la diversité des bouquets, car il} \struck{les pots \ill{}} y aura de la difference dans l'ornement \struck{si on mette l'vn de ces bouquets au milieu ou a l'vn des extremités}.
\note{paragraphe de brouillon, repris plusieurs fois : le problème des 6 pots et des 6 bouquets est corrigé en interligne en un problème de places à pourvoir de commandants, puis tout le bloc est entouré d'un trait et barré de deux longues diagonales. On ne peut dire avec certitude quels mots l'auteur entendait garder ; la lecture ci-dessus suit l'ordre des lignes. Plusieurs mots barrés restent illisibles.}

\uncertain{places d'accord} que 3. on demande en combien de manieres on peut placer 3. de \uncertain{vingt}${}^{+}$

${}^{+}$ Commendans de \uncertain{prix} dans 6. qui sont arrivés, dans 3. de ces places. A

${}^{+}$ On prendra \struck{donc} l'ordre de 3. \uncertain{choses} prises dans 6. multipliant
\note{ces trois lignes, sous le bloc barré, sont liées par des croix de renvoi (+) dont l'ordre de lecture n'est pas sûr ; la dernière s'arrête sur « multipliant », en bas de page, et la suite doit être sur le feuillet suivant. Le « A » final est isolé à droite.}

\page{49}
\note{Page de gauche : verso du feuillet 39, petit feuillet de la même main ronde, collé sur onglet. La page est un brouillon très corrigé : les mots barrés sont nombreux, les corrections écrites au-dessus des lignes, et deux passages sont entourés d'un trait. Le texte continue celui de la vue 48 (« multipliant ») sans titre. Cachet rond « BIBLIOTHÈQUE NATIONALE MANUSCRITS R.F. » à gauche, vers le milieu de la page.}

6. par 5. \& le produit par 4. \& divisant le produit 120. par 6. qui est la combinaison d'ordre de 3. choses, \& on aura 20. pour \struck{\ill{}} le \uncertain{changement} de 3. choses prises dans 6. ou par abreger on multipliera seulement 5. par 4. La combinaison des 3. \add{places prises dans 6. est} \struck{pots dans lesquels on met} pareillement 20. \struck{Les bouquets est pareillement 120. pour y mettre l'ordre} \add{par ce que ces 3 pots \ill{} dans l'ornement}; \struck{\ill{} de la diversité \& que la difference situation des pots} \struck{sont choisis dans 6}. On multipliera donc 120. par 20. \& on aura 2400. façons de \struck{\ill{}} \add{placer 3. Commendans dans chaquun \ill{} de 6. places}.
\note{les lignes 6 à 9 de la page sont reprises trois fois : une première rédaction, sur les « bouquets » et les « pots », est barrée et entourée d'un trait ; les corrections en interligne substituent les places et les commandants. La lecture suit ce qui paraît non barré ; l'ordre exact des morceaux n'est pas sûr. Le « 120 » de « multipliera donc 120 » est surchargé.}

Or il n'importe pas qu'il y ayt autant de \struck{pots} \add{places} que de \struck{bouquets} \add{Commendans}, \add{*} il y en peut avoir plus ou moins. \& la regle sera toujours la mesme.
\note{« 6. places » est écrit en tête de ligne, au-dessus de « Or », et un astérisque renvoie à « il y en peut avoir », en début de la ligne suivante.}

Si par exemple il n'y avoit que 5. \struck{pots} \add{places}, il faudroit prendre la varieté de 3. choses prises dans 5. qui est 60. l'ordre compris \add{qui divisé par 6. on a 10} \& le multipliant par 20. qui est la combinaison des 3. \struck{bouquets} \add{Commendans pris dans 6} on auroit 1200. façons de les \struck{\ill{}} \add{placer}. De mesme s'il y avoit 8. \struck{pots} \add{places} on prendroit la combinaison de 3. pris dans 8. qui est 336. compris l'ordre \add{qui divisé par 6, qui est l'ordre de 3. choses donne 56} qui estant \add{encor} multiplié par 20. donneroit \struck{\ill{}} 1120. façons de \struck{\ill{}} \add{placer} les 3. \struck{bouquets} Commendans.
\note{« 1120 » est écrit au-dessus d'un nombre barré.}

Il y a quelqu'autre chose a considerer dans la combinaison d'ordre, \struck{ou} dans l'assemblage des lettres qui forment les dictions, il y en a quelques vnes qui \struck{\ill{}} \add{\ill{}} les premieres \uncertain{prononcé} quand elles sont ensemble, c'est pourquoy il est necessaire de les separer.

On veut sçavoir par exemple comb. on peut faire de dictions des 8. lettres a, b, c, d, e, i, o, s. a telle condition que les 3. b, c, d. ne se trouvent jamais ensemble. \struck{On considere} Il faudra considerer ces 3. lettres comme vne seule, \& ainsi il n'y aura que 6. choses dont la combinaison est 720. mais par ce que ces 3. lettres se peuvent trouver \add{de suite} \struck{ensemble} en 6. façons, il faut multiplier 720. par 6. le produit est 4320. qu'il faut oster de la combinaison de 8. choses qui est 40320. restera 36000. dictions${}^{+}$ qu'on pourra faire avec ces 8. lettres sans que b, c, d. se trouvent ensemble.
\marginal{${}^{+}$ ou anagrammes}
\note{« ou anagrammes » est écrit dans la marge de gauche, en face de « qu'on pourra faire », et rattaché au mot « dictions » par une croix de renvoi.}

On pourroit encor demander que 2. de ces \uncertain{consones} ne se trouvassent jamais au commencement ni a la fin des dictions. Pour trouver cela il faut voir comb. il \add{se peut faire} \struck{\ill{}} de changemens, supposant que ces 2. de ces lettres sont au commencement ou a la fin; \& par ce qu'il reste 6. lettres outre ces 2. on aura 720. changemens, mais entre ces 720. il y en \struck{aura} a 120. ausquels la 3\textsuperscript{e} consone se trouve contre les 2. autres, \& cela \struck{a esté desja} est compris dans les \uncertain{4320}. qu'il a falu oster de 40320. il faut donc oster 120. de 720. reste 600. qu'il faudra multiplier par 12. par ce que on peut prendre \add{les} 2. lettres dans les 3. en 3. façons, \& a cause de l'ordre il faudra multiplier 3. par 2. \& par ce qu'il se peut faire aussy que les 2. lettres se trouvent a la fin, on aura 12. varietés qui multipliees${}^{\ddagger}$
\marginal{${}^{\ddagger}$ par 600. donnera 7200. qu'il faut oster de 36000. restera 28800 anagrammes}
\note{la fin du calcul est écrite dans la marge de gauche, verticalement, et reliée par le signe $\ddagger$ à la dernière ligne de la page. Le nombre « 4320 » est surchargé : ses deux premiers chiffres sont repassés et la lecture n'est donnée que par le sens.}

\note{Page de droite : recto d'un autre petit feuillet, même main ronde, collé sur onglet ; une pièce nouvelle, avec son titre. En haut à gauche, « N\textsuperscript{o} 2. », à l'encre. En haut à droite, « 40 » de la série continue, et à sa droite « 28 », plus fin, de l'ancienne numérotation. Cachet rond « BIBLIOTHÈQUE NATIONALE MANUSCRITS R.F. » au milieu de la page.}

\subsection*{Hasard}

C'est la coustume a \uncertain{Gennes} d'eslire \add{ou plustost de tirer au sort} tous les ans d'entre les \ill{} cinq personnes a qui \struck{on} donner a avoir les principales charges \struck{des affaires} de la Republique. Cela a donné lieu a des paris qui se font tous les ans touchant ceux a qui le sort arrivera. Il se trouve des Banquiers qui promettront jusques a 20. milles pistoles pour vne qu'on leur donnera si le sort tombe sur 5. qu'on aura nommé; \struck{\ill{}} 5. ou 6. milles s'il n'y a que 4. des 5. qu'on aura nommé, \& 5. ou 600. \struck{pistoles} s'il y en a 3. \struck{\ill{}} pour l'ordinaire ils ne donnent rien pour vn ni pour deux.
\note{« ou plustost de tirer au sort » est écrit au-dessus de la ligne, avec un renvoi. Au début de la deuxième ligne, entre « d'entre les » et « cinq personnes », deux ou trois mots (« Cela … 8 … ») n'ont pu être lus ; ils portent peut-être le nombre de ceux entre qui le sort se tire. Le nom de la ville, dont l'initiale se confond avec un 6, n'est donné que sous réserve.}

On demande quels sont les hazards pour le Banquier, \& pour le pariant, \& quel profit le Banquier peut faire sur \struck{\ill{}} \struck{ce nombre}.

Il faut premierement arrester ce que doit donner le Banquier si \add{le sort tombe sur} ceux qu'on aura nommé, ou sur quelques vns d'eux.

Nous supposons que le Banquier donne 20000. pistoles pour vne si le sort tombe sur les 5. qu'on aura nommé; \struck{\ill{}} 5000. s'il n'y en a que 4. \struck{\ill{}} 300. s'il y en a 3. \& 4. s'il n'y en a que 2.

On verra premierement en combien de manieres les 5. qu'on \add{doit tirer au sort} \struck{nomme} peuvent venir.
\note{« doit tirer au sort » est écrit au-dessus de « nomme » biffé.}

Il faut multiplier 100. par 99. le produit par 98. par 97. \& par 96. mais par ce que le produit de ces 5. nombres contient aussy l'ordre dans lequel ces 5. personnes sont tirées, il le faut diviser par 120.${}^{+}$ ou bien multiplier seulement 80. par 97. 98. \& 99. ce qui est la mesme chose que de multiplier les 5. nombres l'vn par l'autre \& diviser le produit par 120. On aura 75,287,520. qui sont toutes \add{les} façons dont 5. \struck{personnes} \add{billets} peuvent estre tirés \add{\& pris} \struck{sur} \add{dans} 100.
\marginal{${}^{+}$ que est la combinaison de 5 choses}
\note{« billets » est écrit au-dessus de « personnes » biffé. Le nombre est écrit avec des virgules entre les tranches, comme ici.}

Il faut maintenant voir quels sont les hazards du pariant.

Des 5. qu'il a nommé ne pouvant arriver qu'en \add{seule} vne maniere il n'y aura donc qu'vn hazard pour luy \& 75,287,519. pour le Banquier; \& par ce que le Banquier donne 20000. pour 1. il faut \struck{diviser \ill{}} \add{multiplier 1.} par 20000. \struck{on aura 3764 $\frac{3}{4}$ \uncertain{peu plus} qui} \struck{est le profit du Banquier} donc pour \add{20000} \ill{} du hazard qu'a le pariant \struck{le Banquier} le Banquier \struck{\ill{} a 3764} a 75287519. qui divisé par 20000. donnent 3764 $\frac{3}{4}$ \uncertain{et plus}${}^{+}$
\marginal{${}^{+}$ La proportion est donc comme 1. a 3764 $\frac{3}{4}$}
\note{fin de page très corrigée : deux lignes entières sont biffées et une troisième surchargée ; « multiplier 1. » et « 20000 » sont écrits au-dessus de la ligne. La note marginale, à gauche, est rattachée par une croix au résultat « 3764 $\frac{3}{4}$ ».}

Pour avoir les hazards de 4. il faut prendre 5. fois 95. qui
\note{la page s'arrête sur « qui » ; la suite est sur le verso (vue 50).}

\page{50}
\note{Page de gauche : verso du feuillet 40, même main ; le calcul des hasards du pari continue sans rupture (« qui / sont 475. »). L'écriture du recto transparaît dans les marges. En bas à droite, très pâle, « … 1 a 3764 $\frac{3}{4}$ », qui répète la note marginale du recto et paraît en être la transparence.}

sont 475. \struck{pour les hazards de 4} qu'il \add{les} faut oster de tous les hazards ou plustost de ceux qu'a le Banquier sur les hazards de 5. On ostera donc 475. de 75287519. restera 75287044 pour les hazards du Banquier; mais par ce qu'il donne 5000. pour 1. il faut diviser 75287044. par 5000. \& on aura 15057 $\frac{2}{5}$ peu plus pour les hazards du Banquier \& 475. pour ceux du pariant \& divisant l'vn par l'autre, il viendra 31 $\frac{7}{10}$ peu moins pour les hazards du Banquier \& 1. pour ceux du pariant.
\note{« les » est écrit au-dessus de « qu'il faut ». Le 5 de cette main, dans « $\frac{2}{5}$ », a la forme d'un 3.}

Pour avoir les hazards de 3. personnes dans les 5. qu'on a nommé on multipliera par 10 le triangle de 94. qui est 4465. on aura donc 44650 pour les hazards du pariant, qui estant ostés des hazards que le Banquier a eu sur 4. sçavoir de 75287044. restera 75242394. pour les hazards du Banquier qu'il faut diviser par 300. par ce qu'il donne 300. pour 1. \& on aura 250808. peu moins pour les hazards du Banquier, qui estant divisés par 44650. on aura les hazards du Banquier comme 5 $\frac{3}{5}$ peu plus a 1.

Reste a voir les hazards de 2. Pour avoir ceux du pariant on multipliera le tetraedre de 93. sçavoir 138415. par 10. \& on aura 1384150. qu'il faut oster des hazards que le Banquier a eu sur 3. sçavoir de 75242394. restera 73858244. qu'il faut diviser par 4. a cause que le Banquier donne 4. pour 1. \& on aura 18464561. pour les hazards du Banquier, \& les divisant par les hazards du pariant qui sont 1384150. on aura les hazards du Banquier comme 13 $\frac{1}{2}$ peu moins a 1.
\note{la fraction se lit $\frac{1}{2}$ ; la division des deux nombres donnés donne un peu plus de 13 $\frac{1}{3}$. Une tache d'encre précède « 73858244 ».}

On peut aussy considerer les hazards de 1. c'est a dire s'il n'arrive quelqu'vn des 5. qu'on a nommé. Il faudra multiplier par 5. le triangle triangle ou 4\textsuperscript{e}. puissance triangulaire de 92. qui est 3183545. le produit est 15917725. qu'il faut oster des hazards du Banquier sur 2. sçavoir de 73858244. restera 57940519. qui sont les hazards du Banquier; mais par ce qu'il ne donne rien, \struck{s'il} \add{quand} il ne tirera \struck{\ill{}} \add{qu'vn} des 5. qu'on a \struck{\ill{}} nommé, on \struck{\ill{}} divisera \add{le Banquier} 57940519. par 15917725. \& on aura encor 3 $\frac{16}{25}$. de hazard sur 1. qu'aura le pariant, mais ce hazard n'est \add{qu'vn} \struck{\ill{}} profit du Banquier \& le pariant n'y a rien.
\note{« le Banquier » est écrit au-dessus de « divisera » ; « \& le pariant n'y a rien » est serré au début de la ligne suivante.}

On a veu tous ces hazards il les faut assembler. Et premierement \struck{\ill{}} si le sort tombe sur les 5. qui ont esté nommé, le pariant a 20000. S'il en tire 4. il y a 475. hazards pour le pariant, qui multipliés par 5000 que le Banquier doit donner \add{s'il arrive} quelqu'vn des 475. hazards
\note{la page s'arrête ici ; la suite est sur la page de droite.}

\note{Page de droite : recto du feuillet suivant, même main ; la phrase de la page de gauche continue. En haut à droite, « 41 » de la série continue, et à sa droite « 29. », plus fin, de l'ancienne numérotation.}

\struck{on aura} \add{ce sera} 2375000. hazards pour le pariant.

Si le sort tombe sur 3. de ceux qui ont esté nommé, les hazards de 3. sont 44650. qui multipliés par 300. que le Banquier doit donner pour chaqu'vn de ces hazards, \struck{on aura} \add{restera} 13395000. hazards pour le pariant s'il en tire 3. des 5. qu'il a nommé.
\note{« de ces hazards, restera » est écrit dans la marge de gauche et au-dessus de « on aura » biffé.}

Si le sort ne tombe que sur 2. \struck{\ill{}} \add{on a trouvé que les hazards de 2. sont} 1384150. qui multipliés par 4. donnent 5536600. hazards pour le pariant.
\note{la correction « on a trouvé que les hazards de » est écrite au-dessus de la ligne et soulignée ; « sont » est écrit après plusieurs mots biffés.}

\[
\begin{array}{rl}
20000. & \text{hazards de 5.} \\
2375000. & \text{de 4.} \\
13395000. & \text{de 3.} \\
5536600. & \text{de 2.} \\
\hline
21326600. & \text{somme des hazards du pariant}
\end{array}
\]
\note{le tableau est à droite des paragraphes précédents.}

Tous ces hazards assemblés montent a 21326600; \& ce sont les hazards du pariant.

\struck{Tous les hazards assemblés seront 75287520. dont osté 21326600. reste \uncertain{53960920}. qui sont les hazards du Banquier. \uncertain{Les nombres dans les moindres termes sont} comme 533165. a 1349023. qui est comme 1. a 2 $\frac{1}{2}$ ou peu plus.}
\note{ce paragraphe de six lignes est barré de trois grands traits obliques et de traits horizontaux ; plusieurs chiffres en sont incertains. Dans la marge de gauche, en face, une tache d'encre en forme de fanion.}

Pour avoir les hazards du Banquier, il faut \struck{prendre} \add{assembler} tous les hazards du \struck{pari} pariant qui sont 1. 475. 44650. \& 1384150. la somme est 1429276. qui \struck{estant} ostés de tous les hazards qui sont \add{en tout} 75287520. \struck{\ill{}} restera 73858244. pour les hazards du Banquier; les hazards \add{du pariant} seront donc a ceux du Banquier, comme \struck{1429276} \add{21326600} a 73858244. ou dans les moindres termes comme 183850. a 636709. qui est \add{comme 1. a} vn peu moins de \struck{\ill{}} 3 $\frac{1}{2}$ ou justement comme 1. a${}^{+}$
\[
\begin{array}{rl}
1384150. & 2. \\
44650. & 3. \\
475. & 4. \\
1. & \text{de 5} \\
\hline
1429276 &
\end{array}
\]
\marginal{${}^{+}$ 3 $\frac{85159}{183850}$.}
\note{la petite addition est à droite du paragraphe. « 21326600 » est écrit au-dessus de « 1429276 » biffé ; « comme 1. a » au-dessus de « 3 ». La croix renvoie à la marge de gauche, où est écrit « + 3 $\frac{85159}{183850}$ ».}

Ce seroient la les hazards du Banquier \& du pariant si le Banquier ne retiroit rien de ceux a qui le sort est favorable. Mais par ce qu'outre \struck{\ill{}} l'avantage qu'il a dans les hazards il a encor vne pistole de \struck{\ill{}} chaqu'vn de ceux a qui les hazards precedens arrivent, il a pour luy tous les hazards de 5. personnes choisies dans 100. sçavoir 75287520. la proportion \struck{est donc} des hazards \add{du pariant} est donc \add{a ceux du Banquier} comme 21326600. a 75287520. ou comme \struck{\ill{}} 533165. a 1882188. c'est a dire comme 1. a \uncertain{vn} peu plus de 3. \struck{fois} $\frac{1}{2}$.
\note{« fois » est biffé entre « 3. » et « $\frac{1}{2}$ ».}

Mais par ce que d'ordinaire on ne donne rien pour 2. il faut oster \add{des hazards du pariant} les hazards de 2. qui se montent a 5536600. \& il restera 15790000. qui sont aux hazards du Banquier comme 394750. a 1882188. ou comme 1. a \struck{\ill{}} vn peu plus de 4 $\frac{3}{4}$.
\note{après « comme 1. a », un chiffre surchargé et biffé.}

Voicy \struck{\ill{}} le fondement de ces raisons de cette operation.

Pour sçavoir en combien de manieres on peut choisir 5. choses dans 100. on multiplie l'vn par l'autre les 5. nombres 100. 99. 98. 97. 96. \& on divise le dernier produit par l'ordre de 5. choses. Si on ne prend qu'vne chose dans 100. il est evident qu'on ne le pourra faire qu'en 100. façons. Que si on en prend 2. puisque la chose \struck{\ill{}} se prend en 100. façons, avec chaqu'vne \add{des 100.} on peut mettre \struck{\ill{}} laquelle on voudra des 99. restantes, mais on voit que l'ordre \add{y} est \struck{\ill{}} compris par ce que chaqu'vne des 100 sera dans tous les choix la derniere \& la premiere, il faudra donc \struck{\ill{}} diviser par 2. qui sert par l'ordre de 2. choses, le produit de 99. par 100.
\note{la dernière phrase est de lecture difficile par endroits (« sera dans tous les choix la derniere \& la premiere ») ; le reste de la page est blanc. La suite est au verso (vue 51).}

\page{51}
\note{Page de gauche : verso du feuillet 41, même main ; le « fondement » commencé au recto continue. Le bas de la page, blanc, laisse voir à l'envers l'écriture du recto.}

Si on choisit 3. choses dans 100. par ce que 2. choses se prennent en 9900. manieres (c'est le produit de 100. par 99) \& qu'il en reste 98. on pourra choisir \add{chaqu'vne} de ces 98. qui restent \& l'adjouster a chaqu'vne des 9900. façons dont on a choisi 2. choses. Donc le produit de 9900. par 98. contiendra les diverses manieres de choisir 3. choses dans 100. l'ordre compris.
\note{« chaqu'vne » est écrit au-dessus d'un mot biffé.}

Par la mesme raison pour choisir 4 choses il faudra multiplier ce dernier produit par les 97. qui restent; \& pour 5. choses multiplier encor ce dernier produit par 96. \struck{\& \ill{}} \add{on} divisera le produit par 120. qui est l'ordre de 5. choses par ce que par cette construction chaqu'vne des 100. choses tient alternativement chaqu'vn des 5. rangs qui sont dans 5. choses sçavoir le 1\textsuperscript{er} 2\textsuperscript{d} 3\textsuperscript{e} 4\textsuperscript{e} \& dernier.

Pour les hazars du pariant. On multiplie 95. par 5. pour sçavoir en combien de façons il peut tirer 4. \struck{\ill{}} \add{des 5.} qu'il a nommé. Car puisqu'il en manque 1. chaqu'vn des 5. peut manquer, \& en sa place il peut tirer l'vn des 95. \struck{qu'il n'a} dont il n'a nommé aucun, \struck{on} il faut donc multiplier \struck{ce} 95. par 5. \struck{Or} Il est vray que les 4. ostés de 100. il resteroit 96. on ne multiplie pas pourtant par 96. par ce que le 5\textsuperscript{e} \struck{\ill{}} nommé se trouvant au nombre des hazars du pariant, il \struck{ne} seroit \struck{pas} conté deux fois au pariant.

S'il tire 3. des 5. \uncertain{qui ont} esté nommé, les hazars se trouvent en multipliant par 10. le triangle de 94.

On multiplie par 10. par ce que 3. choses se prennent choisies dans 5. en 10 manieres, ainsy qu'il a esté expliqué cy devant, quand on a fait voir combien de façons on peut choisir 5. choses dans 100; car on multipliera 5. 4. 3. l'vn par l'autre \& on divisera le produit 60. par l'ordre de 3. qui est 6.

On \struck{\ill{}} \add{multiplie} par le triangle de 94. par ce que dans les 5. qu'on tire au hazart il n'y \uncertain{a hazart} que 3. \struck{\ill{}} \add{se} trouvent que le pariant a nommé, les 2. autres doivent estre des 95. autres. or dans tous les choix ou hazars le 1\textsuperscript{er} de ces 95. se trouvera avec chaqu'vn des 94. autres. apres le 2\textsuperscript{d} de ces mesmes 95. se trouvera avec chaqu'vn des 93. autres. le 3\textsuperscript{e} avec chaqu'vn des 92. \& ainsy de suite jusques au 94\textsuperscript{e} qui se trouvera avec le dernier des 95. Or ces nombres assemblés font le triangle de 94. par ce qu'il faudroit adjouster 94. avec 93. 92. 91. \&c. jusques a 1.

Pour avoir les hazars de 2. on multiplie \add{par 10} le tetraedre de 93.
\note{la page s'arrête ici, aux deux tiers de sa hauteur, sur une phrase sans suite ; le reste est blanc. La page de droite (feuillet 42) traite d'autre chose.}

\note{Page de droite : un petit feuillet isolé, plus bas que les autres, collé sur onglet ; même main ronde. Deux problèmes nouveaux, sans titre. En haut à droite, « 42 » de la série continue, et au-dessous « 26 », plus fin et plus pâle, de l'ancienne numérotation. Cachet rond « BIBLIOTHÈQUE NATIONALE MANUSCRITS R.F. » à droite, vers le bas.}

On demande a combien de personnes se montent les ancestres en 30. generations, supposant que les mariages ne se soient point faits entre les descendans des premieres \add{\& plus anciennes} generations.

Il faut prendre la 30\textsuperscript{e}. puissance de 2. qui est 1,073,741,824. \& c'est le nombre des ancestres.
\note{un trait horizontal sépare ce problème du suivant.}

Au jeu des eschets, \struck{\ill{}} les \add{8} pions peuvent avancer vne ou deux cases au premier coup. On demande en comb. de façons on peut \add{jouer} ces 8 pions, en ne jouant chaqu'vn d'eux qu'vne seule fois.

S'ils ne pouvoient \add{estre} joués que d'vne seule maniere, on auroit 40,320. façons de les jouer, sçavoir selon la combinaison de l'ordre de 8. choses; mais par ce que chaqu'vn se peut jouer en 2. manieres il faudra multiplier 40320. par la 8\textsuperscript{e}. puissance de 2. sçavoir par 256. \& on aura 10,321,920.

En cette sorte de combinaison ou chaque chose se place en 2. \ill{} manieres il faut multiplier tous les nombres pairs l'vn par l'autre, au lieu qu'en la combinaison simple, on ne multiplie que tous les nombres. Ainsy vne chose se prendra en 2. façons, 2. en 8. 3. en 48. 4. en 384. \&c.
\note{après « en 2. », un signe bref, peut-être une abréviation (« f. ») ; « 8 » est écrit au-dessus d'un chiffre noyé d'encre.}

\[
\begin{array}{rr|l}
2. & 2. & 1 \\
4. & 8 & 2 \\
6. & 48 & 3. \\
8. & 384 & 4 \\
10. & 3840 & 5. \\
12. & 46080 & 6. \\
14. & 645120 & 7. \\
6. & 10,321,920 & 8.
\end{array}
\]
\note{table dans la marge de gauche, en face du paragraphe suivant : à gauche les nombres pairs, au milieu leurs produits successifs, à droite le nombre de pions. Le premier nombre de la dernière ligne se lit « 6. » : le 1 de 16 manque, écrit hors du feuillet ou jamais écrit ; le produit, serré, se lit 10,321,920 comme dans le texte.}

Si chaque chose se pouvoit prendre en 3. manieres, comme si les pions pouvoient avancer vne, deux ou 3. cases; il faudroit multiplier l'vn par l'autre les nombres multiples de 3. sçavoir 3. 6. 9 \&c. \& on auroit 18. pour le changement de 2. choses, 162. pour celuy de 3. \&c. \& ainsy les 8 pions se joueroient la 1\textsuperscript{re} fois en 264,539,520. manieres.

\[
\begin{array}{rr|l}
3. & 3 & 1 \\
6. & 18 & 2 \\
9. & 162 & 3 \\
12. & 1944 & 4 \\
15. & 29160. & 5 \\
18. & 524880 & 6 \\
21. & 11022480 & 7 \\
24. & 264,539,520 & 8
\end{array}
\]
\note{seconde table, à droite du paragraphe. Une tache d'eau touche le paragraphe et la table. Le reste de la page est blanc.}

\end{document}
